\documentclass[11pt]{article}
% Shared print layout. The web's provisional design and MathJax are independent.
\usepackage[utf8]{inputenc}
\usepackage[T1]{fontenc}
\usepackage[spanish,es-nodecimaldot]{babel}
\usepackage{lmodern}
\usepackage{amsmath,amssymb}
\usepackage{graphicx}
\usepackage{geometry}
\geometry{a4paper,left=20mm,right=20mm,top=25mm,bottom=22mm,headheight=16pt,headsep=8mm}
\usepackage{fancyhdr,lastpage}
\usepackage{enumitem}
\usepackage{xcolor}
\usepackage{pgfplots}
\pgfplotsset{compat=1.18}
\usepgfplotslibrary{fillbetween}
\usepackage[hidelinks,pdfencoding=auto]{hyperref}
\hypersetup{pdflang=es-ES,pdfauthor={},pdfsubject={Matemáticas Aplicadas a las Ciencias Sociales II}}
\pagestyle{fancy}
\fancyhf{}
\fancyhead[L]{\small Matemáticas Aplicadas a las Ciencias Sociales II}
\fancyhead[R]{\small 2.\textsuperscript{o} Bachillerato}
\fancyfoot[L]{\footnotesize Programación lineal I | Edición revisada 1.0}
\fancyfoot[R]{\footnotesize Página \thepage\ de \pageref{LastPage}}
\renewcommand{\headrulewidth}{0.4pt}
\setlength{\parindent}{0pt}
\setlength{\parskip}{0.5em}
\setlist{itemsep=0.3em,topsep=0.3em}
\newcommand{\apartado}[1]{\subsection*{#1}}
\pgfplotsset{
  revision/.style={
    width=0.94\linewidth,height=7.5cm,
    axis lines=middle,grid=major,grid style={gray!18},
    tick label style={font=\small},label style={font=\small},
    xlabel={$x$},ylabel={$y$},samples=160,
    legend style={at={(0.5,-0.16)},anchor=north,draw=none,font=\small},
    clip=true
  }
}

\hypersetup{pdftitle={Programación lineal I - Soluciones revisadas}}
\begin{document}
\section*{Programación lineal I: soluciones}
Se sigue el orden del examen: ejercicio 1 obligatorio y \textbf{solo uno} entre 2 y 3. Cada ejercicio vale 5 puntos; la puntuación máxima es 10.
\section*{Ejercicio 1: recinto y mínimo}
La región está definida por
\[x+4y\geq5,\quad x+2y\geq4,\quad7x+5y\leq35,\quad x\geq0.\]
Al intersectar rectas frontera y descartar puntos que incumplen alguna restricción, obtenemos:
\begin{center}
\begin{tabular}{c|c|c}
Vértice & Fronteras que se cortan & $F=2x+y$\\\hline
$(0,2)$ & $x=0$, $x+2y=4$ & $2$\\
$(3,1/2)$ & $x+4y=5$, $x+2y=4$ & $13/2$\\
$(5,0)$ & $x+4y=5$, $7x+5y=35$ & $10$\\
$(0,7)$ & $x=0$, $7x+5y=35$ & $7$
\end{tabular}
\end{center}
Por ejemplo, restando $x+2y=4$ de $x+4y=5$, resulta $y=1/2$ y después $x=3$. Todos los vértices indicados cumplen las cuatro restricciones.

\textbf{Certificado del mínimo para toda la región:}
\[F=\frac12(x+2y)+\frac32x\geq\frac12\cdot4+0=2.\]
La igualdad exige $x=0$ y $x+2y=4$, es decir, $(0,2)$, que es factible. Por tanto,
\[\boxed{F_{\min}=2\quad\text{únicamente en }(0,2)}.\]
\begin{center}
\begin{tikzpicture}
\begin{axis}[revision,height=5.8cm,xmin=-0.5,xmax=5.8,ymin=-0.5,ymax=7.8,xtick={0,1,2,3,4,5},ytick={0,2,4,6,7}]
\addplot[fill=gray!20,draw=black,thick] coordinates {(0,2) (3,0.5) (5,0) (0,7)} -- cycle;
\addplot[densely dotted,domain=0:5] {(5-x)/4};
\addplot[dashed,domain=0:5] {(4-x)/2};
\addplot[domain=0:5] {(35-7*x)/5};
\addplot[only marks,mark=*,mark size=1.7pt] coordinates {(0,2) (3,0.5) (5,0) (0,7)};
\node[anchor=west] at (axis cs:0.15,2.1) {\small mínimo $(0,2)$};
\end{axis}
\end{tikzpicture}
\end{center}
El contorno del polígono está incluido. No se añade $y\geq0$ como hipótesis: se deduce de las restricciones del recinto.
\clearpage
\section*{Ejercicio 2: misiones y programas}
Sean $x$ las misiones e $y$ los programas. Dividiendo los importes por 100\,000 euros, el modelo es
\[
\max F=\frac35x+\frac25y,\qquad
\begin{cases}2x+y\leq24,\\2x+y>20,\\x\geq4,\\y\geq2x.\end{cases}
\]
Las cantidades finales son enteras no negativas. El modelo continuo ayuda a representar y resolver el problema; el óptimo obtenido también es entero.

\textbf{Frontera estricta.} El enunciado dice ``superar'' 2 millones: la desigualdad correcta es $2x+y>20$. Los puntos con inversión exactamente igual a 2 millones quedan excluidos.

La \textbf{clausura} de la región se obtiene sustituyendo temporalmente $>$ por $\geq$ para localizar sus esquinas:
\begin{center}
\begin{tabular}{c|c|c}
Esquina de la clausura & Valor de $F$ & \textquestiondown{}Pertenece a la región?\\\hline
$(4,12)$ & $36/5=7{,}2$ & No: inversión igual a 2 millones\\
$(4,16)$ & $44/5=8{,}8$ & Sí\\
$(6,12)$ & $42/5=8{,}4$ & Sí\\
$(5,10)$ & $7$ & No: inversión igual a 2 millones
\end{tabular}
\end{center}
Se obtienen de intersectar $x=4$, $y=2x$ y las dos rectas de inversión. Por ejemplo, con $x=4$ y $2x+y=24$, resulta $y=16$.

\textbf{No basta con evaluar esquinas excluidas.} Para garantizar un máximo alcanzado en la región original, comprobamos para cualquier punto factible:
\[
F=\frac25(2x+y)-\frac15x
\leq\frac25\cdot24-\frac15\cdot4=\frac{44}5.
\]
La igualdad exige $2x+y=24$ y $x=4$, por lo que $y=16$.
\[\boxed{x=4,\qquad y=16,\qquad F_{\max}=8{,}8}.\]
\textbf{Comprobación:} la inversión es $200\,000\cdot4+100\,000\cdot16=2\,400\,000$ euros, superior a 2 millones y dentro del presupuesto. Hay al menos 4 misiones y $4\leq16/2$.
\clearpage
\section*{Ejercicio 2: representación y cantidades enteras}
\begin{center}
\begin{tikzpicture}
\begin{axis}[revision,height=8cm,xmin=3.5,xmax=6.8,ymin=8.5,ymax=17.5,xtick={4,5,6},ytick={10,12,14,16},xlabel={Misiones ($x$)},ylabel={Programas ($y$)}]
\addplot[fill=gray!20,draw=none] coordinates {(4,12) (4,16) (6,12) (5,10)} -- cycle;
\addplot[thick] coordinates {(4,12) (4,16) (6,12) (5,10)};
\addplot[dashed,thick] coordinates {(4,12) (5,10)};
\addplot[only marks,mark=o,mark size=3pt,mark options={fill=white,solid}] coordinates {(4,12) (5,10)};
\addplot[only marks,mark=*,mark size=2.3pt] coordinates {(4,16) (6,12)};
\node[anchor=south west] at (axis cs:4,16) {\small óptimo $(4,16)$};
\node[anchor=north west] at (axis cs:5,10) {\small $(5,10)$ excluido};
\node[anchor=north east] at (axis cs:4,12) {\small $(4,12)$};
\node[anchor=west] at (axis cs:6,12) {\small $(6,12)$};
\end{axis}
\end{tikzpicture}
\end{center}
El segmento discontinuo $2x+y=20$ \textbf{no pertenece} a la región; sus dos extremos se marcan con círculos vacíos. Las otras aristas están incluidas, salvo esos extremos.

Para cantidades enteras, $2x+y>20$ equivale a $2x+y\geq21$. De $y\geq2x$ y $2x+y\leq24$ se deduce $x\leq6$; junto a $x\geq4$, solo hay tres valores posibles:
\begin{center}
\begin{tabular}{c|c}
Misiones & Programas enteros factibles\\\hline
$4$ & $13,14,15,16$\\
$5$ & $11,12,13,14$\\
$6$ & $12$
\end{tabular}
\end{center}
Son nueve posibilidades. El punto $(4,16)$ es entero y alcanza la cota superior ya demostrada; también es el óptimo del problema con cantidades enteras.
\clearpage
\section*{Ejercicio 3: baterías y beneficio máximo}
Sean $x$ las baterías A e $y$ las baterías B. El modelo es
\[
\max Z=130x+140y,\qquad
\begin{cases}x+y\geq10,\\y-x\leq10,\\3x+2y\leq120,\\x\geq0,\quad y\geq0.\end{cases}
\]
La tercera restricción procede de dividir $150x+100y\leq6\,000$ entre 50. Los beneficios del enunciado ya son beneficios unitarios: no se restan de nuevo los gastos de producción.

Las intersecciones factibles y sus valores son
\begin{center}
\begin{tabular}{c|r}
Vértice & $Z$ (euros)\\\hline
$(0,10)$ & $1\,400$\\
$(10,0)$ & $1\,300$\\
$(40,0)$ & $5\,200$\\
$(20,30)$ & $6\,800$
\end{tabular}
\end{center}
La intersección de $y=x+10$ y $3x+2y=120$ da $5x=100$, luego $(x,y)=(20,30)$.

\textbf{Certificado para toda la región:}
\[Z=54(3x+2y)+32(y-x)\leq54\cdot120+32\cdot10=6\,800.\]
En $(20,30)$ se alcanza la igualdad y se cumplen $x+y=50\geq10$, $y-x=10$ y $150x+100y=6\,000$. Así,
\[\boxed{20\text{ baterías A},\quad30\text{ baterías B},\quad Z_{\max}=6\,800\text{ euros}}.\]
El óptimo es entero y único, pues la igualdad exige simultáneamente las dos rectas anteriores.
\begin{center}
\begin{tikzpicture}
\begin{axis}[revision,height=5.2cm,xmin=-2,xmax=44,ymin=-2,ymax=36,xtick={0,10,20,30,40},ytick={0,10,20,30},xlabel={Baterías A ($x$)},ylabel={Baterías B ($y$)}]
\addplot[fill=gray!20,draw=black,thick] coordinates {(0,10) (10,0) (40,0) (20,30)} -- cycle;
\addplot[only marks,mark=*,mark size=2pt] coordinates {(0,10) (10,0) (40,0) (20,30)};
\node[anchor=south] at (axis cs:20,30) {\small óptimo $(20,30)$};
\end{axis}
\end{tikzpicture}
\end{center}
\end{document}
